%%%%%%%%%%%%%%%%%%%%%%%%%%%%%%%%%%%%%%%%%%%%%%%%%%%%%%%%%%%%%
%% CATEGORY THEORY — Chapters 9–11 + End Matter
%% Monoidal Categories / Grothendieck Toposes / ∞-Categories
%%%%%%%%%%%%%%%%%%%%%%%%%%%%%%%%%%%%%%%%%%%%%%%%%%%%%%%%%%%%%

%%%%%%%%%%%%%%%%%%%%%%%%%%%%%%%%%%%%%%%%%%%%%%%%%%%%%%%%%%%%
%%         CHAPTER 9: MONOIDAL AND BRAIDED CATEGORIES     %%
%%%%%%%%%%%%%%%%%%%%%%%%%%%%%%%%%%%%%%%%%%%%%%%%%%%%%%%%%%%%

\chapter{Monoidal and Braided Categories}\label{ch:monoidal}
\minitoc

\vspace{1em}

In many areas of mathematics and theoretical physics, one encounters
categories equipped with a ``tensor product''---a bifunctor that combines
objects in a way that is associative and unital up to coherent isomorphism.
Such categories are called \emph{monoidal}.  They provide the natural
setting for studying algebras, coalgebras, and their modules in a
category-theoretic framework.  When the tensor product is commutative
up to isomorphism, one obtains \emph{braided} and \emph{symmetric} monoidal
categories, which underpin the theory of quantum groups and topological
quantum field theories.  In this chapter we develop the foundations of
monoidal category theory, state Mac~Lane's coherence theorem, study
enriched and closed categories, and touch upon the Curry--Howard--Lambek
correspondence.

\index{monoidal category|(}

%% ============================================================
\section{Monoidal categories}\label{sec:monoidal-cats}
%% ============================================================

\begin{definition}[Monoidal category]\label{def:monoidal-cat}
\index{monoidal category!definition}
\index{tensor product!in a monoidal category}
\index{associator}
\index{unitor!left}
\index{unitor!right}
\index{unit object}
A \textbf{monoidal category} is a tuple
$(\mathcal{C},\, \otimes,\, I,\, \alpha,\, \lambda,\, \rho)$ where:
\begin{enumerate}[label=(\roman*)]
\item $\mathcal{C}$ is a category;
\item $\otimes \colon \mathcal{C}\times\mathcal{C}\to\mathcal{C}$ is a
      bifunctor called the \textbf{tensor product};
\item $I \in \Ob(\mathcal{C})$ is an object called the \textbf{unit object}
      (or \textbf{monoidal unit});
\item $\alpha_{A,B,C}\colon (A\otimes B)\otimes C
      \xrightarrow{\;\sim\;} A\otimes(B\otimes C)$
      is a natural isomorphism called the \textbf{associator};
\item $\lambda_A \colon I\otimes A \xrightarrow{\;\sim\;} A$
      is a natural isomorphism called the \textbf{left unitor};
\item $\rho_A \colon A\otimes I \xrightarrow{\;\sim\;} A$
      is a natural isomorphism called the \textbf{right unitor};
\end{enumerate}
subject to the \emph{pentagon axiom} and the \emph{triangle axiom} below.
\end{definition}

\begin{definition}[Pentagon axiom]\label{def:pentagon}
\index{pentagon axiom}
For all objects $A,B,C,D$ of~$\mathcal{C}$, the following diagram commutes:
\[
\begin{tikzcd}[column sep=small]
& (A\otimes B)\otimes(C\otimes D) \arrow[dr, "\alpha_{A,B,C\otimes D}"] & \\
((A\otimes B)\otimes C)\otimes D
  \arrow[ur, "\alpha_{A\otimes B,C,D}"]
  \arrow[d, "\alpha_{A,B,C}\otimes\id_D"'] & &
A\otimes(B\otimes(C\otimes D)) \\
(A\otimes(B\otimes C))\otimes D
  \arrow[rr, "\alpha_{A,B\otimes C,D}"'] & &
A\otimes((B\otimes C)\otimes D)
  \arrow[u, "\id_A\otimes\alpha_{B,C,D}"']
\end{tikzcd}
\]
\end{definition}

\begin{definition}[Triangle axiom]\label{def:triangle}
\index{triangle axiom}
For all objects $A,B$ of~$\mathcal{C}$, the following diagram commutes:
\[
\begin{tikzcd}
(A\otimes I)\otimes B
  \arrow[rr, "\alpha_{A,I,B}"]
  \arrow[dr, "\rho_A\otimes\id_B"'] & &
A\otimes(I\otimes B)
  \arrow[dl, "\id_A\otimes\lambda_B"] \\
& A\otimes B &
\end{tikzcd}
\]
\end{definition}

\begin{definition}[Strict monoidal category]\label{def:strict-monoidal}
\index{monoidal category!strict}
A monoidal category is called \textbf{strict} if $\alpha$, $\lambda$,
and $\rho$ are all identity natural transformations, i.e.\
$(A\otimes B)\otimes C = A\otimes (B\otimes C)$,
$I\otimes A = A$, and $A\otimes I = A$ on the nose.
\end{definition}

\begin{example}\label{ex:monoidal-examples}
\index{monoidal category!examples}
The following are important examples of monoidal categories.
\begin{enumerate}[label=(\roman*)]
\item $(\Set, \times, \{*\})$: the category of sets with cartesian product.
      \index{Set@$\Set$!monoidal structure}
\item $(\Ab, \otimes_\Z, \Z)$: abelian groups with the tensor product over~$\Z$.
      \index{Ab@$\Ab$!monoidal structure}
\item $(\Vect_\K, \otimes_\K, \K)$: vector spaces over a field~$\K$
      with the usual tensor product.
      \index{Vect@$\Vect$!monoidal structure}
\item $(R\text{-}\Mod, \otimes_R, R)$ for a commutative ring~$R$.
      \index{Mod@$\Mod$!monoidal structure}
\item $(\Cat, \times, \mathbf{1})$: the category of small categories
      with the cartesian product.
      \index{Cat@$\Cat$!monoidal structure}
\item $(\mathrm{End}(\mathcal{C}), \circ, \id_\mathcal{C})$: endofunctors
      of a category~$\mathcal{C}$ with composition.  This is a strict
      monoidal category.
      \index{endofunctor!monoidal category of}
\item $(\Set, \sqcup, \varnothing)$: sets with disjoint union.
\item The category of chain complexes $\mathrm{Ch}(R)$ with the
      tensor product of complexes.
      \index{chain complex!monoidal structure}
\end{enumerate}
\end{example}

\begin{remark}\label{rem:monoidal-functor-notation}
We often write just $(\mathcal{C}, \otimes, I)$ for a monoidal category,
leaving the structural isomorphisms implicit.
\end{remark}


%% ============================================================
\section{Mac Lane's coherence theorem}\label{sec:coherence}
%% ============================================================

\index{coherence theorem|(}
\index{Mac Lane!coherence theorem}

The pentagon and triangle axioms may appear insufficient to guarantee that
\emph{all} diagrams built from $\alpha$, $\lambda$, $\rho$ and their
inverses commute.  Mac~Lane's celebrated coherence theorem asserts that
this is indeed the case.

\begin{theorem}[Mac Lane's coherence theorem]\label{thm:maclane-coherence}
\index{coherence theorem!Mac Lane}
Every monoidal category is monoidally equivalent to a strict monoidal
category.  Equivalently, every diagram in a monoidal category whose
morphisms are composites of instances of $\alpha$, $\alpha^{-1}$,
$\lambda$, $\lambda^{-1}$, $\rho$, $\rho^{-1}$, identities, and tensor
products thereof, commutes.
\end{theorem}

\begin{proof}[Proof sketch]
One constructs a strict monoidal category $\mathcal{C}_s$ whose objects
are finite words in $\Ob(\mathcal{C})$ (including the empty word, which
serves as the strict unit) and whose morphisms are induced by those
of~$\mathcal{C}$.  The multiplication is word concatenation.
An explicit monoidal equivalence $\mathcal{C}_s \simeq \mathcal{C}$ is
built by iterating the tensor product.  The commutativity of arbitrary
constraint diagrams then reduces to the trivially commuting diagrams
in~$\mathcal{C}_s$.  For the full proof, see Mac~Lane~\cite{MacLane1963}
or~\cite{MacLaneCWM}.
\end{proof}

\begin{remark}\label{rem:coherence-practical}
\index{coherence theorem!practical consequence}
The practical consequence of the coherence theorem is that, for
calculations, one may pretend that every monoidal category is strict:
one simply omits all parentheses and all instances of the
associator and unitors.  We shall make frequent use of this
simplification in what follows.
\end{remark}

\index{coherence theorem|)}


%% ============================================================
\section{Braided and symmetric monoidal categories}\label{sec:braided}
%% ============================================================

\index{braided monoidal category|(}

\begin{definition}[Braided monoidal category]\label{def:braided}
\index{braided monoidal category!definition}
\index{braiding}
A \textbf{braided monoidal category} is a monoidal category
$(\mathcal{C}, \otimes, I)$ equipped with a natural isomorphism
\[
   \sigma_{A,B} \colon A\otimes B \xrightarrow{\;\sim\;} B\otimes A
\]
called the \textbf{braiding}, satisfying the two \emph{hexagon axioms}:
\[
\begin{tikzcd}[column sep=1.8em]
(A\otimes B)\otimes C
  \arrow[r, "\alpha"]
  \arrow[d, "\sigma_{A,B}\otimes\id"'] &
A\otimes(B\otimes C)
  \arrow[r, "\sigma_{A,B\otimes C}"] &
(B\otimes C)\otimes A
  \arrow[d, "\alpha"] \\
(B\otimes A)\otimes C
  \arrow[r, "\alpha"'] &
B\otimes(A\otimes C)
  \arrow[r, "\id\otimes\sigma_{A,C}"'] &
B\otimes(C\otimes A)
\end{tikzcd}
\]
and
\[
\begin{tikzcd}[column sep=1.8em]
A\otimes(B\otimes C)
  \arrow[r, "\alpha^{-1}"]
  \arrow[d, "\id\otimes\sigma_{B,C}"'] &
(A\otimes B)\otimes C
  \arrow[r, "\sigma_{A\otimes B,C}"] &
C\otimes(A\otimes B)
  \arrow[d, "\alpha^{-1}"] \\
A\otimes(C\otimes B)
  \arrow[r, "\alpha^{-1}"'] &
(A\otimes C)\otimes B
  \arrow[r, "\sigma_{A,C}\otimes\id"'] &
(C\otimes A)\otimes B
\end{tikzcd}
\]
\end{definition}

\begin{definition}[Symmetric monoidal category]\label{def:symmetric}
\index{symmetric monoidal category}
\index{braiding!symmetric}
A braided monoidal category is \textbf{symmetric} if
$\sigma_{B,A}\circ \sigma_{A,B} = \id_{A\otimes B}$ for all $A,B$.
\end{definition}

\begin{example}\label{ex:braided-symmetric}
\begin{enumerate}[label=(\roman*)]
\item $(\Set, \times, \{*\})$, $(\Ab, \otimes_\Z, \Z)$, and
      $(\Vect_\K, \otimes_\K, \K)$ are all symmetric monoidal categories
      with the evident swap maps.
      \index{symmetric monoidal category!examples}
\item The category of $R$-modules with $\otimes_R$ is symmetric monoidal.
\item The category of graded $R$-modules with the Koszul sign rule
      $\sigma(a\otimes b) = (-1)^{|a||b|}\, b\otimes a$ is symmetric monoidal.
      \index{Koszul sign rule}
\item The category of representations of a quantum group is braided
      monoidal but typically not symmetric.
      \index{quantum group}
\end{enumerate}
\end{example}

\index{braided monoidal category|)}


%% ============================================================
\section{Monoid objects}\label{sec:monoid-objects}
%% ============================================================

\index{monoid object|(}

\begin{definition}[Monoid object]\label{def:monoid-object}
\index{monoid object!definition}
\index{multiplication morphism}
\index{unit morphism}
Let $(\mathcal{C}, \otimes, I)$ be a monoidal category.
A \textbf{monoid object} (or \textbf{monoid}) in~$\mathcal{C}$ is a
triple $(M, \mu, \eta)$ where $M\in\Ob(\mathcal{C})$,
$\mu\colon M\otimes M \to M$ is the \textbf{multiplication}, and
$\eta\colon I\to M$ is the \textbf{unit}, satisfying associativity
and unitality:
\[
\begin{tikzcd}
(M\otimes M)\otimes M \arrow[r, "\alpha"] \arrow[d, "\mu\otimes\id"'] &
M\otimes(M\otimes M) \arrow[r, "\id\otimes\mu"] &
M\otimes M \arrow[d, "\mu"] \\
M\otimes M \arrow[rr, "\mu"'] && M
\end{tikzcd}
\qquad
\begin{tikzcd}
I\otimes M \arrow[r, "\eta\otimes\id"] \arrow[dr, "\lambda"'] &
M\otimes M \arrow[d, "\mu"] &
M\otimes I \arrow[l, "\id\otimes\eta"'] \arrow[dl, "\rho"] \\
& M &
\end{tikzcd}
\]
A \textbf{morphism of monoid objects} $f\colon (M,\mu,\eta)\to (M',\mu',\eta')$
is a morphism $f\colon M\to M'$ in~$\mathcal{C}$ with
$f\circ\mu = \mu'\circ(f\otimes f)$ and $f\circ\eta = \eta'$.
The category of monoid objects in~$\mathcal{C}$ is denoted
$\mathrm{Mon}(\mathcal{C})$.
\index{Mon@$\mathrm{Mon}(\mathcal{C})$}
\end{definition}

\begin{example}\label{ex:monoid-objects}
\index{monoid object!examples}
\begin{enumerate}[label=(\roman*)]
\item A monoid object in $(\Set, \times, \{*\})$ is a monoid in the
      usual sense.
\item A monoid object in $(\Ab, \otimes_\Z, \Z)$ is a ring (not
      necessarily commutative).
      \index{ring!as monoid object}
\item A monoid object in $(\Vect_\K, \otimes_\K, \K)$ is a
      $\K$-algebra.
      \index{algebra!as monoid object}
\item A monoid object in $(\mathrm{End}(\mathcal{C}), \circ, \id_\mathcal{C})$
      is a monad on~$\mathcal{C}$.
      \index{monad!as monoid object}
\item A \textbf{commutative monoid object} in a symmetric monoidal
      category $(\mathcal{C}, \otimes, I, \sigma)$ is a monoid $(M,\mu,\eta)$
      with $\mu \circ \sigma_{M,M} = \mu$.
      \index{monoid object!commutative}
\end{enumerate}
\end{example}

\begin{definition}[Comonoid object]\label{def:comonoid}
\index{comonoid object}
Dually, a \textbf{comonoid object} in $(\mathcal{C}, \otimes, I)$ is
a triple $(C, \Delta, \varepsilon)$ where
$\Delta\colon C\to C\otimes C$ (the \emph{comultiplication}) and
$\varepsilon\colon C\to I$ (the \emph{counit}) satisfy the duals of
the associativity and unitality diagrams.
\end{definition}

\index{monoid object|)}


%% ============================================================
\section{Enriched categories}\label{sec:enriched}
%% ============================================================

\index{enriched category|(}

Ordinary category theory is ``enriched over~$\Set$'': for each pair of
objects $A,B$, the hom $\Hom(A,B)$ is a \emph{set}.  But in practice
these hom-sets often carry additional structure---they may be abelian
groups ($\Ab$-enriched), topological spaces ($\Top$-enriched), or
chain complexes (dg-enriched).

\begin{definition}[$\mathcal{V}$-enriched category]\label{def:enriched-cat}
\index{enriched category!definition}
\index{V-category@$\mathcal{V}$-category}
Let $(\mathcal{V}, \otimes, I)$ be a monoidal category.
A \textbf{$\mathcal{V}$-enriched category} (or \textbf{$\mathcal{V}$-category})
$\mathcal{A}$ consists of:
\begin{enumerate}[label=(\roman*)]
\item a class $\Ob(\mathcal{A})$ of objects;
\item for each pair $A,B\in\Ob(\mathcal{A})$, an object
      $\mathcal{A}(A,B)\in\Ob(\mathcal{V})$ (the \emph{hom-object});
\item for each triple $A,B,C$, a \emph{composition morphism}
      \[
          c_{A,B,C}\colon \mathcal{A}(B,C)\otimes\mathcal{A}(A,B)
          \longrightarrow \mathcal{A}(A,C) \quad\text{in }\mathcal{V};
      \]
\item for each object~$A$, an \emph{identity morphism}
      $j_A\colon I\to\mathcal{A}(A,A)$ in~$\mathcal{V}$;
\end{enumerate}
subject to associativity and unitality axioms expressed as commutative
diagrams in~$\mathcal{V}$.
\end{definition}

\begin{example}\label{ex:enriched-examples}
\index{enriched category!examples}
\begin{enumerate}[label=(\roman*)]
\item A $\Set$-enriched category is an ordinary (locally small) category.
\item An $\Ab$-enriched category is a \emph{preadditive category}:
      hom-sets are abelian groups and composition is bilinear.
      \index{preadditive category}
\item A $\Top$-enriched category is a \emph{topological category}:
      hom-sets are topological spaces and composition is continuous.
\item A category enriched over the monoidal category $([0,\infty]^{\op}, +, 0)$
      (extended non-negative reals, reversed order, addition) is a
      (generalised) \emph{metric space} in the sense of Lawvere.
      \index{Lawvere metric space}
      \index{metric space!as enriched category}
\item A category enriched over chain complexes is a \emph{dg-category}.
      \index{dg-category}
\end{enumerate}
\end{example}

\begin{definition}[$\mathcal{V}$-functor]\label{def:enriched-functor}
\index{enriched category!functor}
\index{V-functor@$\mathcal{V}$-functor}
A \textbf{$\mathcal{V}$-functor} $F\colon\mathcal{A}\to\mathcal{B}$
between $\mathcal{V}$-categories consists of a map
$F\colon\Ob(\mathcal{A})\to\Ob(\mathcal{B})$ and morphisms
$F_{A,B}\colon\mathcal{A}(A,B)\to\mathcal{B}(FA,FB)$ in~$\mathcal{V}$
compatible with composition and identities.
\end{definition}

\begin{definition}[$\mathcal{V}$-natural transformation]\label{def:enriched-nat}
\index{enriched category!natural transformation}
A \textbf{$\mathcal{V}$-natural transformation}
$\eta\colon F\Rightarrow G$ between $\mathcal{V}$-functors
$F,G\colon\mathcal{A}\to\mathcal{B}$ consists of a family of morphisms
$\eta_A\colon I\to\mathcal{B}(FA,GA)$ in~$\mathcal{V}$, one for each
$A\in\Ob(\mathcal{A})$, satisfying the enriched naturality condition.
\end{definition}

\index{enriched category|)}


%% ============================================================
\section{Closed monoidal and cartesian closed categories}%
\label{sec:closed-monoidal}
%% ============================================================

\index{closed monoidal category|(}
\index{cartesian closed category|(}

\begin{definition}[Closed monoidal category]\label{def:closed-monoidal}
\index{closed monoidal category!definition}
\index{internal hom}
A monoidal category $(\mathcal{C}, \otimes, I)$ is
\textbf{(left) closed} if for each object~$B$ the functor
$(-)\otimes B\colon\mathcal{C}\to\mathcal{C}$ has a right adjoint,
denoted $[B, -]$ or $\underline{\Hom}(B,-)$:
\[
    \Hom_\mathcal{C}(A\otimes B,\, C)
    \;\cong\;
    \Hom_\mathcal{C}(A,\, [B,C])
    \quad\text{naturally in } A \text{ and } C.
\]
The object $[B,C]$ is called the \textbf{internal hom}
(or \textbf{exponential object}).
\end{definition}

\begin{definition}[Cartesian closed category]\label{def:cartesian-closed}
\index{cartesian closed category!definition}
\index{exponential object}
A category $\mathcal{C}$ with finite products is \textbf{cartesian closed}
if it is closed monoidal with respect to the cartesian monoidal structure
$(\mathcal{C}, \times, 1)$.  That is, for every object~$B$ the functor
$(-)\times B$ has a right adjoint $(-)^B$:
\[
    \Hom(A\times B,\, C) \;\cong\; \Hom(A,\, C^B).
\]
\end{definition}

\begin{example}\label{ex:cartesian-closed}
\index{cartesian closed category!examples}
\begin{enumerate}[label=(\roman*)]
\item $\Set$ is cartesian closed with $C^B = \Hom_{\Set}(B,C)$.
      \index{Set@$\Set$!cartesian closed}
\item The category $\Cat$ of small categories is cartesian closed with
      the exponential $\mathcal{D}^\mathcal{C} = \Fun(\mathcal{C},\mathcal{D})$.
\item Any elementary topos is cartesian closed (Theorem~\ref{thm:topos-ccc}).
\item $\Top$ is \emph{not} cartesian closed in general; one must pass to
      a convenient subcategory such as compactly generated Hausdorff spaces.
      \index{Top@$\Top$!not cartesian closed}
\item $\Ab$ is not cartesian closed (the cartesian product is the direct
      sum, and $\Hom_\Ab(A\oplus B, C)\not\cong\Hom_\Ab(A, C^B)$ in general),
      but it is closed monoidal with $\otimes_\Z$ and internal hom $\Hom_\Z(B,C)$.
\end{enumerate}
\end{example}

\begin{proposition}[Evaluation and coevaluation]\label{prop:eval-coeval}
\index{evaluation morphism}
\index{coevaluation morphism}
In a closed monoidal category, there are natural morphisms:
\begin{enumerate}[label=(\alph*)]
\item the \textbf{evaluation}: $\mathrm{ev}_{B,C}\colon [B,C]\otimes B\to C$,
      corresponding to $\id_{[B,C]}$ under the adjunction;
\item the \textbf{coevaluation}: $\mathrm{coev}_{A,B}\colon A\to [B, A\otimes B]$,
      corresponding to $\id_{A\otimes B}$ under the adjunction.
\end{enumerate}
\end{proposition}

\begin{proof}
Apply the adjunction isomorphism
$\Hom(A\otimes B, C)\cong\Hom(A,[B,C])$
to the appropriate identity morphisms.
\end{proof}

\index{closed monoidal category|)}
\index{cartesian closed category|)}


%% ============================================================
\section{The Curry--Howard--Lambek correspondence}%
\label{sec:curry-howard-lambek}
%% ============================================================

\index{Curry--Howard--Lambek correspondence|(}
\index{lambda calculus}
\index{intuitionistic logic}

The Curry--Howard--Lambek correspondence is a remarkable three-way
dictionary between logic, computation, and category theory.  We give a
brief account aimed at illustrating the categorical perspective.

\begin{center}
\renewcommand{\arraystretch}{1.3}
\begin{tabular}{lll}
\hline
\textbf{Logic} & \textbf{Type theory} & \textbf{Category theory} \\
\hline
Proposition $A$ & Type $A$ & Object $A$ \\
Proof of $A$ & Term of type $A$ & Morphism $1\to A$ \\
$A\Rightarrow B$ & Function type $A\to B$ & Exponential $B^A$ \\
$A\wedge B$ & Product type $A\times B$ & Product $A\times B$ \\
$A\vee B$ & Sum type $A + B$ & Coproduct $A + B$ \\
$\top$ (true) & Unit type $1$ & Terminal object \\
$\bot$ (false) & Empty type $0$ & Initial object \\
Modus ponens & Function application & Evaluation $B^A\times A\to B$ \\
\hline
\end{tabular}
\end{center}

\begin{theorem}[Lambek]\label{thm:lambek}
\index{Lambek!correspondence}
\index{cartesian closed category!and lambda calculus}
The simply-typed lambda calculus (with products) is the internal language
of cartesian closed categories.  Specifically:
\begin{enumerate}[label=(\alph*)]
\item Every cartesian closed category gives rise to a simply-typed
      lambda calculus whose types are the objects and whose terms are
      the morphisms.
\item Every simply-typed lambda calculus (modulo $\beta\eta$-equivalence)
      gives rise to a cartesian closed category, the \emph{syntactic category}.
\item These two constructions are mutually inverse (up to equivalence).
\end{enumerate}
\end{theorem}

\begin{remark}\label{rem:curry-howard-extensions}
\index{Curry--Howard--Lambek correspondence!extensions}
The correspondence extends far beyond the simply-typed case:
\begin{itemize}
\item Dependent types correspond to locally cartesian closed categories.
\item Linear logic corresponds to $*$-autonomous categories (a form of
      closed monoidal category).
      \index{linear logic}
      \index{*-autonomous category@$*$-autonomous category}
\item Homotopy type theory corresponds to $(\infty,1)$-toposes.
      \index{homotopy type theory}
\end{itemize}
\end{remark}

\index{Curry--Howard--Lambek correspondence|)}


%% ============================================================
\section{Exercises}\label{sec:monoidal-exercises}
%% ============================================================

\begin{exercise}\label{ex:strict-is-monoidal}
Show that every strict monoidal category is a monoidal category
(the axioms are trivially satisfied with identity natural transformations).
\end{exercise}

\begin{exercise}\label{ex:monoidal-unique-unit}
\index{unit object!uniqueness}
Let $(\mathcal{C}, \otimes, I)$ be a monoidal category.  Show that the
unit object~$I$ is unique up to isomorphism, in the sense that if
$(I', \lambda', \rho')$ also satisfies the axioms, then $I\cong I'$.
\end{exercise}

\begin{exercise}\label{ex:commutative-monoid-object}
Let $(\mathcal{C}, \otimes, I, \sigma)$ be a symmetric monoidal category.
Show that the category $\mathrm{CMon}(\mathcal{C})$ of commutative monoid
objects in~$\mathcal{C}$ admits a monoidal structure inherited from that
of~$\mathcal{C}$.
\end{exercise}

\begin{exercise}\label{ex:comonoid-in-set}
\index{comonoid object!in $\Set$}
Show that every set $S$ has a unique comonoid structure in
$(\Set, \times, \{*\})$, given by $\Delta(s) = (s,s)$ and
$\varepsilon(s) = *$.
\end{exercise}

\begin{exercise}\label{ex:enriched-underlying}
\index{underlying category}
Let $\mathcal{A}$ be a $\mathcal{V}$-enriched category where
$\mathcal{V}$ is closed monoidal.  Define the \emph{underlying ordinary
category} $\mathcal{A}_0$ of~$\mathcal{A}$ and show it is indeed a category.
\textit{Hint:} take $\Hom_{\mathcal{A}_0}(A,B) = \Hom_\mathcal{V}(I, \mathcal{A}(A,B))$.
\end{exercise}

\begin{exercise}\label{ex:day-convolution}
\index{Day convolution}
Let $(\mathcal{C}, \otimes, I)$ be a small monoidal category.
Show that the category of presheaves $[\mathcal{C}^{\op}, \Set]$ carries
a monoidal structure given by \emph{Day convolution}:
\[
    (F\star G)(C)
    = \int^{A,B\in\mathcal{C}} \Hom_\mathcal{C}(A\otimes B, C)
      \times F(A)\times G(B).
\]
Identify the unit object of this monoidal structure.
\end{exercise}

\begin{exercise}\label{ex:ccc-lambda}
Verify explicitly that in $\Set$, the adjunction
$\Hom(A\times B, C)\cong\Hom(A, C^B)$ corresponds to currying of functions:
$f\colon A\times B\to C$ is sent to $\hat{f}\colon A\to C^B$ defined by
$\hat{f}(a)(b) = f(a,b)$.
\end{exercise}

\index{monoidal category|)}


%%%%%%%%%%%%%%%%%%%%%%%%%%%%%%%%%%%%%%%%%%%%%%%%%%%%%%%%%%%%
%%    CHAPTER 10: GROTHENDIECK TOPOSES — INTRODUCTION     %%
%%%%%%%%%%%%%%%%%%%%%%%%%%%%%%%%%%%%%%%%%%%%%%%%%%%%%%%%%%%%

\chapter{Grothendieck Toposes---Introduction}\label{ch:toposes}
\minitoc

\vspace{1em}

Topos theory, introduced by Grothendieck and his school in the 1960s
for the purposes of algebraic geometry, has grown into one of the deepest
and most far-reaching branches of category theory.  A \emph{Grothendieck
topos} is a category of sheaves on a site; it simultaneously generalises
the category of sheaves on a topological space and the category of
presheaves on a small category.  Elementary toposes, introduced by
Lawvere and Tierney, isolate the key categorical properties that make
a category ``behave like the category of sets.''

In this chapter we give an introduction to the subject, defining sites,
Grothendieck topologies, sheaves, and the resulting toposes.
We discuss elementary toposes, subobject classifiers, and the internal
logic, then state Giraud's theorem characterising Grothendieck toposes.

\index{topos|(}
\index{Grothendieck topos|(}

%% ============================================================
\section{Sites and Grothendieck topologies}\label{sec:sites}
%% ============================================================

\index{site|(}
\index{Grothendieck topology|(}

\begin{definition}[Sieve]\label{def:sieve}
\index{sieve}
Let $\mathcal{C}$ be a category and $C\in\Ob(\mathcal{C})$.
A \textbf{sieve on~$C$} is a collection~$S$ of morphisms with
codomain~$C$ such that if $f\colon D\to C$ is in~$S$ and
$g\colon E\to D$ is any morphism, then $f\circ g\in S$.
Equivalently, a sieve on~$C$ is a subfunctor of the representable
presheaf $\Hom(-,C)$.
\end{definition}

\begin{definition}[Grothendieck topology]\label{def:grothendieck-topology}
\index{Grothendieck topology!definition}
\index{covering sieve}
A \textbf{Grothendieck topology}~$J$ on a small category~$\mathcal{C}$
assigns to each object~$C$ a collection $J(C)$ of sieves on~$C$, called
\textbf{covering sieves}, satisfying:
\begin{enumerate}[label=(\roman*)]
\item \textbf{Maximal sieve:} The maximal sieve
      $t_C = \{f \mid \mathrm{cod}(f) = C\}$ is in $J(C)$.
\item \textbf{Stability:} If $S\in J(C)$ and $h\colon D\to C$ is any
      morphism, then the pullback sieve
      $h^*S = \{g\colon E\to D \mid h\circ g\in S\}$ is in $J(D)$.
\item \textbf{Transitivity:} If $S\in J(C)$ and $R$ is a sieve on~$C$
      such that for every $f\colon D\to C$ in~$S$ the pullback
      $f^*R\in J(D)$, then $R\in J(C)$.
\end{enumerate}
\end{definition}

\begin{definition}[Site]\label{def:site}
\index{site!definition}
A \textbf{site} is a pair $(\mathcal{C}, J)$ consisting of a small
category~$\mathcal{C}$ and a Grothendieck topology~$J$ on~$\mathcal{C}$.
\end{definition}

\begin{definition}[Coverage]\label{def:coverage}
\index{coverage}
\index{covering family}
In practice, Grothendieck topologies are often specified via
\textbf{coverages}: for each $C\in\Ob(\mathcal{C})$, one gives a
collection of families $\{f_i\colon C_i\to C\}_{i\in I}$, called
\textbf{covering families}, satisfying a stability condition under
pullback.  The covering families generate sieves, hence a Grothendieck
topology.
\end{definition}

\begin{example}\label{ex:topologies}
\index{Grothendieck topology!examples}
\begin{enumerate}[label=(\roman*)]
\item \textbf{Trivial topology.}
      \index{trivial topology}
      Only the maximal sieves cover.  Then every presheaf is a sheaf.
\item \textbf{Discrete topology.}
      \index{discrete topology}
      Every sieve is a covering sieve.  The only sheaf is the terminal
      presheaf.
\item \textbf{Canonical topology on $\Top$.}
      \index{canonical topology}
      For an open set~$U$, a covering sieve is generated by an
      open cover $\{U_i\}$ of~$U$.
\item \textbf{Zariski topology.}
      \index{Zariski topology}
      On the category of commutative rings (with the opposite orientation),
      covering families correspond to collections of localisations
      $\{R\to R[f_i^{-1}]\}$ where $(f_1,\ldots,f_n) = R$.
\item \textbf{\'Etale topology.}
      \index{etale topology@\'etale topology}
      On the category of schemes, covering families are jointly surjective
      families of \'etale morphisms.
\item \textbf{fppf and fpqc topologies.}
      \index{fppf topology}
      \index{fpqc topology}
      Finer topologies on schemes used in algebraic geometry.
\end{enumerate}
\end{example}

\index{Grothendieck topology|)}
\index{site|)}


%% ============================================================
\section{Sheaves on a site}\label{sec:sheaves-site}
%% ============================================================

\index{sheaf|(}

\begin{definition}[Presheaf]\label{def:presheaf-reminder}
\index{presheaf}
A \textbf{presheaf} on a small category~$\mathcal{C}$ is a functor
$F\colon\mathcal{C}^{\op}\to\Set$.  The category of presheaves is
$\mathrm{PSh}(\mathcal{C}) = [\mathcal{C}^{\op}, \Set]$.
\index{PSh@$\mathrm{PSh}(\mathcal{C})$}
\end{definition}

\begin{definition}[Sheaf on a site]\label{def:sheaf-site}
\index{sheaf!on a site}
\index{sheaf condition}
Let $(\mathcal{C}, J)$ be a site.  A presheaf $F\colon\mathcal{C}^{\op}\to\Set$
is a \textbf{sheaf} (for the topology~$J$) if for every covering sieve
$S\in J(C)$, the natural map
\[
    F(C) \longrightarrow \varprojlim_{(f\colon D\to C)\in S} F(D)
\]
is a bijection.  Equivalently, for every covering family
$\{f_i\colon C_i\to C\}$ generating a covering sieve, the diagram
\[
\begin{tikzcd}
F(C) \arrow[r] &
\displaystyle\prod_i F(C_i)
  \arrow[r, shift left=0.5ex]
  \arrow[r, shift right=0.5ex] &
\displaystyle\prod_{i,j} F(C_i\times_C C_j)
\end{tikzcd}
\]
is an equaliser.
\end{definition}

\begin{definition}[Category of sheaves]\label{def:sheaf-cat}
\index{Sh@$\mathrm{Sh}(\mathcal{C}, J)$}
\index{topos!Grothendieck}
The full subcategory of $\mathrm{PSh}(\mathcal{C})$ consisting of
sheaves for the topology~$J$ is denoted
$\mathrm{Sh}(\mathcal{C}, J)$.  It is called a
\textbf{Grothendieck topos}.
\end{definition}

\begin{theorem}[Sheafification]\label{thm:sheafification}
\index{sheafification}
\index{associated sheaf}
The inclusion $i\colon\mathrm{Sh}(\mathcal{C},J)\hookrightarrow
\mathrm{PSh}(\mathcal{C})$ has a left adjoint
$a\colon\mathrm{PSh}(\mathcal{C})\to\mathrm{Sh}(\mathcal{C},J)$
called \textbf{sheafification} (or the \textbf{associated sheaf functor}).
Moreover:
\begin{enumerate}[label=(\alph*)]
\item $a$ is left exact (preserves finite limits);
\item $a\circ i \cong \id$;
\item the adjunction $a\dashv i$ is a reflection.
\end{enumerate}
\end{theorem}

\begin{proof}[Proof sketch]
The sheafification $aF$ is constructed by the ``plus construction''
applied twice: $(F^+)^+$.  Given $F$, one defines
\[
    F^+(C) = \varinjlim_{S\in J(C)} \mathrm{Match}(S, F)
\]
where $\mathrm{Match}(S,F)$ is the set of matching families for the
sieve~$S$.  One application of $(-)^+$ makes~$F$ separated; a second
application yields a sheaf.  Left exactness follows from the fact that
the colimit over covering sieves is filtered.
\end{proof}

\begin{example}\label{ex:sheaves-on-space}
\index{sheaf!on a topological space}
If $X$ is a topological space and $\mathcal{O}(X)$ is the poset of
open subsets (viewed as a category), then the open-cover topology makes
$(\mathcal{O}(X), J)$ a site and $\mathrm{Sh}(\mathcal{O}(X), J)$ is
the classical category of sheaves on~$X$.
\end{example}

\index{sheaf|)}


%% ============================================================
\section{Elementary toposes}\label{sec:elementary-topos}
%% ============================================================

\index{elementary topos|(}

\begin{definition}[Elementary topos]\label{def:elementary-topos}
\index{elementary topos!definition}
\index{subobject classifier}
An \textbf{elementary topos} is a category $\mathcal{E}$ satisfying:
\begin{enumerate}[label=(\roman*)]
\item $\mathcal{E}$ has all finite limits;
\item $\mathcal{E}$ has all finite colimits;
\item $\mathcal{E}$ is cartesian closed;
\item $\mathcal{E}$ has a \textbf{subobject classifier}.
\end{enumerate}
\end{definition}

\begin{definition}[Subobject classifier]\label{def:subobject-classifier}
\index{subobject classifier!definition}
\index{Omega@$\Omega$}
\index{true morphism}
A \textbf{subobject classifier} in a category~$\mathcal{E}$ with
finite limits is a monomorphism $\mathrm{true}\colon 1\to\Omega$
such that for every monomorphism $m\colon S\rightarrowtail A$
there exists a unique morphism $\chi_m\colon A\to\Omega$ making the
following square a pullback:
\[
\begin{tikzcd}
S \arrow[r] \arrow[d, tail, "m"'] & 1 \arrow[d, "\mathrm{true}"] \\
A \arrow[r, "\chi_m"'] & \Omega
\end{tikzcd}
\]
The morphism $\chi_m$ is called the \textbf{characteristic morphism}
(or \textbf{classifying morphism}) of the subobject~$m$.
\index{characteristic morphism}
\index{classifying morphism}
\end{definition}

\begin{example}\label{ex:subobject-classifier}
\index{subobject classifier!examples}
\begin{enumerate}[label=(\roman*)]
\item In $\Set$, $\Omega = \{0,1\}$ with $\mathrm{true}(*)=1$.
      The classifying morphism of a subset $S\subseteq A$ is
      the characteristic function $\chi_S$.
      \index{Set@$\Set$!subobject classifier}
\item In the topos of sheaves $\mathrm{Sh}(X)$ on a topological
      space~$X$, $\Omega$ is the sheaf of open subsets:
      $\Omega(U) = \{V\subseteq U \mid V\text{ open}\}$.
\item In a presheaf topos $[\mathcal{C}^{\op}, \Set]$,
      $\Omega(C) = \{\text{sieves on } C\}$.
      \index{presheaf topos!subobject classifier}
\end{enumerate}
\end{example}

\begin{theorem}\label{thm:topos-ccc}
\index{elementary topos!is cartesian closed}
Every elementary topos is cartesian closed.  In particular, every
Grothendieck topos is cartesian closed.
\end{theorem}

\begin{proof}[Proof sketch]
The power object $P(B) = \Omega^B$ exists by assumption in an elementary
topos (cartesian closedness is part of the definition).  For a
Grothendieck topos, cartesian closedness follows from the fact that
$\mathrm{Sh}(\mathcal{C},J)$ is a reflective subcategory of the
cartesian closed presheaf category, and the reflection preserves
finite products.
\end{proof}

\begin{proposition}\label{prop:topos-properties}
\index{elementary topos!properties}
Every elementary topos $\mathcal{E}$ satisfies:
\begin{enumerate}[label=(\alph*)]
\item $\mathcal{E}$ is balanced (every monic epic is an isomorphism);
\item Epimorphisms in $\mathcal{E}$ are coequalizers;
\item Every morphism $f\colon A\to B$ factors as a regular epimorphism
      followed by a monomorphism (image factorisation).
\end{enumerate}
\end{proposition}

\index{elementary topos|)}


%% ============================================================
\section{Internal logic}\label{sec:internal-logic}
%% ============================================================

\index{internal logic|(}
\index{Mitchell--B\'enabou language}

Every topos has an \emph{internal logic} that allows one to reason about
objects and morphisms using a language that resembles set-theoretic
reasoning, but is \emph{intuitionistic}: the law of excluded middle need
not hold.

\begin{definition}[Internal logic of a topos]\label{def:internal-logic}
\index{internal logic!definition}
The \textbf{internal logic} (or \textbf{Mitchell--B\'enabou language})
of a topos~$\mathcal{E}$ interprets:
\begin{itemize}
\item \emph{types} as objects of~$\mathcal{E}$;
\item \emph{terms} of type~$A$ in context~$\Gamma$ as morphisms
      $\Gamma\to A$;
\item \emph{propositions} in context~$\Gamma$ as morphisms $\Gamma\to\Omega$
      (subobjects of~$\Gamma$);
\item logical connectives via operations on~$\Omega$:
      $\wedge$, $\vee$, $\Rightarrow$, $\neg$, $\forall$, $\exists$.
\end{itemize}
\end{definition}

\begin{remark}\label{rem:internal-logic-intuitionistic}
\index{intuitionistic logic!in a topos}
\index{law of excluded middle}
\index{Boolean topos}
The internal logic of a topos is \textbf{intuitionistic higher-order logic}.
The law of excluded middle $\phi\vee\neg\phi$ holds internally if and only
if $\Omega\cong 1\sqcup 1$, in which case $\mathcal{E}$ is called a
\textbf{Boolean topos}.  The topos $\Set$ is Boolean; sheaf toposes on
non-trivial spaces are typically non-Boolean.
\end{remark}

\begin{remark}\label{rem:kripke-joyal}
\index{Kripke--Joyal semantics}
The precise semantics for interpreting formulae of the internal language
in a topos is called \textbf{Kripke--Joyal semantics} (or \emph{forcing}).
It provides the tool for transferring set-theoretic proofs to arbitrary
toposes.
\end{remark}

\index{internal logic|)}


%% ============================================================
\section{Geometric morphisms}\label{sec:geometric-morphisms}
%% ============================================================

\index{geometric morphism|(}

\begin{definition}[Geometric morphism]\label{def:geometric-morphism}
\index{geometric morphism!definition}
\index{direct image functor}
\index{inverse image functor}
A \textbf{geometric morphism} between toposes
$f\colon\mathcal{E}\to\mathcal{F}$ is an adjunction $f^*\dashv f_*$
where:
\begin{itemize}
\item $f^*\colon\mathcal{F}\to\mathcal{E}$ is called the \textbf{inverse
      image functor};
\item $f_*\colon\mathcal{E}\to\mathcal{F}$ is called the \textbf{direct
      image functor};
\item $f^*$ is required to be left exact (preserve finite limits).
\end{itemize}
\end{definition}

\begin{example}\label{ex:geometric-morphisms}
\index{geometric morphism!examples}
\begin{enumerate}[label=(\roman*)]
\item A continuous map $f\colon X\to Y$ of topological spaces induces
      a geometric morphism $f\colon\mathrm{Sh}(X)\to\mathrm{Sh}(Y)$
      with $f_*$ the direct image of sheaves and $f^*$ the inverse
      image.
\item For any topos $\mathcal{E}$, there is an essentially unique
      geometric morphism $\gamma\colon\mathcal{E}\to\Set$ given by
      $\gamma^* = \Delta$ (constant sheaf) and $\gamma_* = \Gamma$
      (global sections).
      \index{global sections}
      \index{constant sheaf functor}
\item The sheafification adjunction $a\dashv i$ for a site $(\mathcal{C},J)$
      is a geometric morphism
      $\mathrm{Sh}(\mathcal{C},J)\to\mathrm{PSh}(\mathcal{C})$.
\end{enumerate}
\end{example}

\begin{definition}[Categories of geometric morphisms]\label{def:geom-morph-cats}
\index{geometric morphism!natural transformation}
A \textbf{2-cell} between geometric morphisms $f,g\colon\mathcal{E}\to\mathcal{F}$
is a natural transformation $\alpha\colon f^*\Rightarrow g^*$ (equivalently,
$\beta\colon g_*\Rightarrow f_*$ by the mate correspondence).  This gives
a 2-category $\mathfrak{Top}$ of (Grothendieck) toposes.
\index{Top (2-category of toposes)@$\mathfrak{Top}$}
\end{definition}

\index{geometric morphism|)}


%% ============================================================
\section{Giraud's theorem}\label{sec:giraud}
%% ============================================================

\index{Giraud's theorem|(}

Giraud's theorem gives a purely categorical characterisation of
Grothendieck toposes, without reference to sites.

\begin{theorem}[Giraud]\label{thm:giraud}
\index{Giraud's theorem!statement}
A category $\mathcal{E}$ is a Grothendieck topos (i.e.\ equivalent to
$\mathrm{Sh}(\mathcal{C},J)$ for some site $(\mathcal{C},J)$)
if and only if it satisfies the following conditions:
\begin{enumerate}[label=(\roman*)]
\item $\mathcal{E}$ has all small colimits;
\item colimits in $\mathcal{E}$ are universal (stable under pullback);
      \index{universal colimit}
\item coproducts in $\mathcal{E}$ are disjoint;
      \index{disjoint coproduct}
\item every equivalence relation in $\mathcal{E}$ is effective
      (is the kernel pair of its coequaliser);
      \index{effective equivalence relation}
\item $\mathcal{E}$ has a small set of generators.
      \index{generator!small set of}
\end{enumerate}
\end{theorem}

\begin{remark}\label{rem:giraud-consequences}
\index{Giraud's theorem!consequences}
Giraud's theorem implies, in particular, that every Grothendieck topos:
\begin{itemize}
\item is an elementary topos;
\item has all small limits and colimits;
\item is a locally presentable category;
\item satisfies the axiom of choice internally if and only if every
      epimorphism splits.
      \index{axiom of choice!in a topos}
\end{itemize}
\end{remark}

\index{Giraud's theorem|)}


%% ============================================================
\section{Exercises}\label{sec:topos-exercises}
%% ============================================================

\begin{exercise}\label{ex:sieves-subfunctors}
Show that sieves on an object~$C$ in a category~$\mathcal{C}$ are in
bijection with subfunctors of the representable presheaf $\Hom(-,C)$.
\end{exercise}

\begin{exercise}\label{ex:presheaf-topos}
\index{presheaf topos}
Let $\mathcal{C}$ be a small category.  Show that $[\mathcal{C}^{\op}, \Set]$
is a Grothendieck topos (for the trivial topology).
Identify the subobject classifier explicitly.
\end{exercise}

\begin{exercise}\label{ex:omega-internal-heyting}
\index{Heyting algebra}
\index{subobject classifier!Heyting algebra structure}
Let $\mathcal{E}$ be an elementary topos.  Show that the subobject
classifier $\Omega$ carries the structure of an \emph{internal Heyting
algebra}: there exist morphisms
$\wedge, \vee, \Rightarrow\colon\Omega\times\Omega\to\Omega$
satisfying the Heyting algebra axioms internally.
\end{exercise}

\begin{exercise}\label{ex:geometric-morphism-compose}
Show that geometric morphisms compose: if
$f\colon\mathcal{E}\to\mathcal{F}$ and $g\colon\mathcal{F}\to\mathcal{G}$
are geometric morphisms, then so is $g\circ f$ with
$(g\circ f)^* = f^*\circ g^*$ and $(g\circ f)_* = g_*\circ f_*$.
\end{exercise}

\begin{exercise}\label{ex:sheaf-condition-equaliser}
\index{sheaf condition!equaliser form}
Let $(\mathcal{C},J)$ be a site and $\{f_i\colon C_i\to C\}_{i\in I}$
a covering family.  Show that a presheaf~$F$ satisfies the sheaf condition
for this covering if and only if the diagram
\[
\begin{tikzcd}
F(C) \arrow[r] &
\displaystyle\prod_{i\in I} F(C_i)
  \arrow[r, shift left=0.5ex, "d_0"]
  \arrow[r, shift right=0.5ex, "d_1"'] &
\displaystyle\prod_{(i,j)\in I\times I} F(C_i\times_C C_j)
\end{tikzcd}
\]
is an equaliser, where $d_0$ and $d_1$ are induced by the two projections
from the fibre product.
\end{exercise}

\begin{exercise}\label{ex:boolean-topos}
\index{Boolean topos}
Show that a topos $\mathcal{E}$ is Boolean if and only if every subobject
has a complement, i.e.\ for every mono $m\colon S\rightarrowtail A$
there exists $m'\colon S'\rightarrowtail A$ with $S\sqcup S'\cong A$
and $S\cap S'\cong 0$.
\end{exercise}

\begin{exercise}\label{ex:etale-site}
\index{etale topology@\'etale topology}
Describe the \'etale site of $\mathrm{Spec}(\Z)$ and explain why
the resulting topos is related to the absolute Galois group
$\mathrm{Gal}(\overline{\Q}/\Q)$.
\end{exercise}

\index{Grothendieck topos|)}
\index{topos|)}


%%%%%%%%%%%%%%%%%%%%%%%%%%%%%%%%%%%%%%%%%%%%%%%%%%%%%%%%%%%%
%%   CHAPTER 11: ∞-CATEGORIES — OVERVIEW                  %%
%%%%%%%%%%%%%%%%%%%%%%%%%%%%%%%%%%%%%%%%%%%%%%%%%%%%%%%%%%%%

\chapter{$\infty$-Categories---Overview}\label{ch:infinity-cats}
\minitoc

\vspace{1em}

Classical category theory considers objects, morphisms between objects,
and equalities between morphisms.  But in many contexts---homotopy theory,
homological algebra, derived algebraic geometry---one needs morphisms
between morphisms (2-cells), morphisms between those (3-cells), and so on,
ad infinitum, with all compositions associative and unital only up to
coherent higher-dimensional cells.  This is the domain of
\emph{higher category theory} and, in particular,
\emph{$(\infty,1)$-categories}.

This chapter provides a survey of the main ideas, definitions, and
results.  Proofs are mostly omitted; we aim to give the reader a map
of the landscape and precise pointers to the literature.

\index{infinity-category@$\infty$-category|(}
\index{higher category theory|(}

%% ============================================================
\section{Motivation}\label{sec:infty-motivation}
%% ============================================================

\index{infinity-category@$\infty$-category!motivation}

The need for higher categories arises from several independent sources.

\begin{enumerate}[label=(\roman*)]
\item \textbf{Homotopy theory.}
\index{homotopy theory}
The homotopy category $\mathrm{Ho}(\Top)$ of topological spaces
loses too much information: mapping spaces are truncated to sets.
One wants a ``category'' where $\Hom(X,Y)$ is a space (or a
simplicial set), composition is associative up to homotopy, and
all higher coherences are recorded.

\item \textbf{Derived categories.}
\index{derived category}
The derived category $D(\mathcal{A})$ of an abelian category~$\mathcal{A}$
is obtained by formally inverting quasi-isomorphisms, but
the resulting category is poorly behaved (e.g.\ it lacks functorial
cones).  The $\infty$-categorical enhancement---the
\emph{derived $\infty$-category}---remedies these defects.

\item \textbf{Stacks and higher stacks.}
\index{stack!higher}
In algebraic geometry, stacks are ``sheaves of groupoids.''
Higher stacks are sheaves of $\infty$-groupoids, and their natural
home is the $\infty$-topos.

\item \textbf{Topological quantum field theory.}
\index{topological quantum field theory}
The cobordism hypothesis of Baez--Dolan (proved by Lurie) requires
the language of $(\infty,n)$-categories.
\end{enumerate}


%% ============================================================
\section{Quasi-categories}\label{sec:quasi-categories}
%% ============================================================

\index{quasi-category|(}
\index{Joyal}
\index{Boardman--Vogt}

The most developed model for $(\infty,1)$-categories is that of
\emph{quasi-categories}, introduced by Boardman--Vogt and extensively
developed by Joyal and Lurie.

\begin{definition}[Simplicial set]\label{def:simplicial-set-reminder}
\index{simplicial set}
\index{simplex category}
Let $\Delta$ denote the simplex category (finite non-empty ordinals and
order-preserving maps).  A \textbf{simplicial set} is a functor
$X\colon\Delta^{\op}\to\Set$.  We write $X_n = X([n])$ for the set of
$n$-simplices.  The category of simplicial sets is
$\mathrm{sSet} = [\Delta^{\op}, \Set]$.
\index{sSet@$\mathrm{sSet}$}
\end{definition}

\begin{definition}[Horn]\label{def:horn}
\index{horn}
\index{horn!inner}
\index{horn!outer}
For $0\le k\le n$, the \textbf{$k$-th horn} $\Lambda^n_k\subset\Delta^n$
is the simplicial subset obtained by removing the interior of~$\Delta^n$
and the face opposite vertex~$k$.  The horn is called \textbf{inner} if
$0 < k < n$ and \textbf{outer} otherwise.
\end{definition}

\begin{definition}[Quasi-category]\label{def:quasi-category}
\index{quasi-category!definition}
\index{inner horn!filler}
\index{weak Kan complex}
A simplicial set $X$ is a \textbf{quasi-category} (or \textbf{$\infty$-category},
or \textbf{weak Kan complex}) if every inner horn has a filler:
for every $0 < k < n$ and every map $\Lambda^n_k\to X$, there exists
an extension $\Delta^n\to X$:
\[
\begin{tikzcd}
\Lambda^n_k \arrow[r] \arrow[d, hook] & X \\
\Delta^n \arrow[ur, dashed] &
\end{tikzcd}
\]
Note that the filler is \emph{not required to be unique}.
\end{definition}

\begin{remark}\label{rem:kan-vs-quasi}
\index{Kan complex}
\index{infinity-groupoid@$\infty$-groupoid}
If one requires fillers for \emph{all} horns (inner and outer), one
obtains a \textbf{Kan complex}, which models an \emph{$\infty$-groupoid}:
every 1-simplex is invertible up to higher homotopy.  A quasi-category
is to an $\infty$-category what a Kan complex is to an $\infty$-groupoid.
\end{remark}

\begin{example}\label{ex:quasi-cat-examples}
\index{quasi-category!examples}
\begin{enumerate}[label=(\roman*)]
\item The nerve $N(\mathcal{C})$ of an ordinary category~$\mathcal{C}$ is
      a quasi-category in which inner horn fillers are \emph{unique}.
      \index{nerve!as quasi-category}
\item The singular simplicial set $\mathrm{Sing}(X)$ of a topological
      space~$X$ is a Kan complex, hence a quasi-category.
      \index{singular simplicial set}
\item The \textbf{dg-nerve} of a dg-category is a quasi-category.
      \index{dg-nerve}
\item The \textbf{homotopy coherent nerve} of a simplicial category is a
      quasi-category.
      \index{homotopy coherent nerve}
\end{enumerate}
\end{example}

\begin{definition}[Homotopy category of a quasi-category]\label{def:homotopy-cat-qcat}
\index{homotopy category!of a quasi-category}
Given a quasi-category~$X$, its \textbf{homotopy category} $\mathrm{h}X$
is the ordinary category with:
\begin{itemize}
\item objects: the vertices ($0$-simplices) of~$X$;
\item morphisms $x\to y$: homotopy classes of edges ($1$-simplices)
      from~$x$ to~$y$, where two edges are homotopic if they are
      related by a $2$-simplex with degenerate third edge;
\item composition: given by choosing fillers for inner $2$-horns.
\end{itemize}
\end{definition}

\index{quasi-category|)}


%% ============================================================
\section{Models for $(\infty,1)$-categories}\label{sec:models-infty1}
%% ============================================================

\index{(infinity,1)-category@$(\infty,1)$-category|(}

There are several equivalent models for $(\infty,1)$-categories.

\begin{definition}[$(\infty,1)$-category---informal]\label{def:infty1-informal}
\index{(infinity,1)-category@$(\infty,1)$-category!informal definition}
An \textbf{$(\infty,1)$-category} is a ``category enriched in spaces''---a
structure with objects, morphisms, homotopies between morphisms,
homotopies between homotopies, etc., where all $k$-morphisms for $k\ge 2$
are invertible (up to higher morphisms).
\end{definition}

\begin{remark}\label{rem:models-infty1}
\index{(infinity,1)-category@$(\infty,1)$-category!models}
The principal models, all Quillen equivalent, include:
\begin{enumerate}[label=(\roman*)]
\item \textbf{Quasi-categories} (simplicial sets with inner horn fillers)
      --- Joyal, Lurie.
      \index{quasi-category}
\item \textbf{Complete Segal spaces} --- Rezk.
      \index{complete Segal space}
      \index{Rezk}
\item \textbf{Segal categories} --- Hirschowitz--Simpson.
      \index{Segal category}
\item \textbf{Simplicial categories} ($\mathrm{sSet}$-enriched categories)
      --- Bergner.
      \index{simplicial category}
      \index{Bergner}
\item \textbf{Relative categories} (categories with weak equivalences)
      --- Barwick--Kan.
      \index{relative category}
\end{enumerate}
The fact that these models are all equivalent is a deep theorem
(the ``comparison theorem''), proved using Quillen model structures.
\index{comparison theorem}
\end{remark}

\index{(infinity,1)-category@$(\infty,1)$-category|)}


%% ============================================================
\section{The $\infty$-Yoneda lemma}\label{sec:infty-yoneda}
%% ============================================================

\index{Yoneda lemma!infinity version@$\infty$-version|(}

The Yoneda lemma generalises to the $\infty$-categorical setting.

\begin{definition}[$\infty$-presheaf]\label{def:infty-presheaf}
\index{infinity-presheaf@$\infty$-presheaf}
Let $\mathcal{C}$ be a small $(\infty,1)$-category.
The \textbf{$\infty$-category of presheaves} on~$\mathcal{C}$ is
$\mathcal{P}(\mathcal{C}) = \Fun(\mathcal{C}^{\op}, \mathcal{S})$,
where $\mathcal{S}$ denotes the $\infty$-category of spaces
(i.e.\ $\infty$-groupoids, modelled by Kan complexes).
\index{S@$\mathcal{S}$!infinity-category of spaces@$\infty$-category of spaces}
\end{definition}

\begin{theorem}[$\infty$-Yoneda lemma]\label{thm:infty-yoneda}
\index{Yoneda lemma!infinity version@$\infty$-version}
\index{Yoneda embedding!infinity version@$\infty$-version}
Let $\mathcal{C}$ be a small $(\infty,1)$-category.  The Yoneda embedding
\[
    \mathbf{y}\colon \mathcal{C}\hookrightarrow\mathcal{P}(\mathcal{C}),
    \quad C\mapsto\mathrm{Map}_\mathcal{C}(-,C)
\]
is fully faithful.  Moreover, for any presheaf
$F\in\mathcal{P}(\mathcal{C})$ and any $C\in\mathcal{C}$,
\[
    \mathrm{Map}_{\mathcal{P}(\mathcal{C})}(\mathbf{y}(C), F)
    \;\simeq\; F(C)
\]
as objects of~$\mathcal{S}$ (an equivalence of spaces).
\end{theorem}

\begin{remark}\label{rem:infty-yoneda-universal}
\index{free cocompletion!infinity-categorical@$\infty$-categorical}
The $\infty$-category $\mathcal{P}(\mathcal{C})$ is the
\textbf{free cocompletion} of~$\mathcal{C}$ under small colimits,
just as in the 1-categorical case.  Every $\infty$-presheaf is a
colimit of representables.
\end{remark}

\index{Yoneda lemma!infinity version@$\infty$-version|)}


%% ============================================================
\section{$\infty$-adjunctions and $\infty$-limits}\label{sec:infty-adj}
%% ============================================================

\index{infinity-adjunction@$\infty$-adjunction|(}

\begin{definition}[$\infty$-adjunction]\label{def:infty-adjunction}
\index{infinity-adjunction@$\infty$-adjunction!definition}
An \textbf{$\infty$-adjunction} between $(\infty,1)$-categories
$\mathcal{C}$ and $\mathcal{D}$ is a pair of functors
$F\colon\mathcal{C}\to\mathcal{D}$, $G\colon\mathcal{D}\to\mathcal{C}$
together with an equivalence of mapping spaces
\[
    \mathrm{Map}_\mathcal{D}(FC, D)
    \;\simeq\;
    \mathrm{Map}_\mathcal{C}(C, GD)
\]
natural in~$C$ and~$D$.
\end{definition}

\begin{remark}\label{rem:infty-adjunction-properties}
\index{infinity-adjunction@$\infty$-adjunction!properties}
The theory of $\infty$-adjunctions parallels the 1-categorical theory:
\begin{enumerate}[label=(\roman*)]
\item $\infty$-adjoints are unique up to contractible choice;
\item an $\infty$-left adjoint preserves $\infty$-colimits;
\item an $\infty$-right adjoint preserves $\infty$-limits;
\item the $\infty$-adjoint functor theorem holds (under appropriate
      presentability hypotheses).
      \index{adjoint functor theorem!infinity version@$\infty$-version}
\end{enumerate}
\end{remark}

\begin{definition}[$\infty$-limit]\label{def:infty-limit}
\index{infinity-limit@$\infty$-limit}
\index{homotopy limit}
Let $F\colon\mathcal{J}\to\mathcal{C}$ be a diagram in an
$(\infty,1)$-category~$\mathcal{C}$.  An \textbf{$\infty$-limit}
of~$F$ is an object $L\in\mathcal{C}$ together with equivalences
\[
    \mathrm{Map}_\mathcal{C}(X, L)
    \;\simeq\;
    \lim_{j\in\mathcal{J}}\, \mathrm{Map}_\mathcal{C}(X, Fj)
\]
natural in~$X$, where the right-hand side is a homotopy limit of spaces.
Dually for $\infty$-colimits.
\end{definition}

\index{infinity-adjunction@$\infty$-adjunction|)}


%% ============================================================
\section{Higher toposes}\label{sec:higher-toposes}
%% ============================================================

\index{infinity-topos@$\infty$-topos|(}

\begin{definition}[$\infty$-topos]\label{def:infty-topos}
\index{infinity-topos@$\infty$-topos!definition}
An \textbf{$\infty$-topos} is an $(\infty,1)$-category that is an
accessible left exact localisation of an $\infty$-presheaf category
$\mathcal{P}(\mathcal{C})$ for some small $(\infty,1)$-category~$\mathcal{C}$.
\end{definition}

\begin{remark}\label{rem:infty-topos-giraud}
\index{infinity-topos@$\infty$-topos!Giraud-type characterisation}
Lurie proves an $\infty$-categorical analogue of Giraud's theorem:
an $(\infty,1)$-category $\mathcal{E}$ is an $\infty$-topos if and only if:
\begin{enumerate}[label=(\roman*)]
\item $\mathcal{E}$ is presentable (cocomplete and accessible);
\item colimits in $\mathcal{E}$ are universal;
\item groupoid objects in $\mathcal{E}$ are effective.
\end{enumerate}
Note that the disjointness condition for coproducts, present in Giraud's
classical theorem, is automatic in the $\infty$-setting.
\end{remark}

\begin{example}\label{ex:infty-toposes}
\index{infinity-topos@$\infty$-topos!examples}
\begin{enumerate}[label=(\roman*)]
\item The $\infty$-category $\mathcal{S}$ of spaces (Kan complexes) is
      the initial $\infty$-topos, playing the role of~$\Set$.
\item For a topological space~$X$, the $\infty$-category
      $\mathrm{Sh}_\infty(X)$ of $\infty$-sheaves (sheaves of spaces)
      on~$X$ is an $\infty$-topos.
      \index{infinity-sheaf@$\infty$-sheaf}
\item Derived algebraic geometry uses $\infty$-toposes as the ambient
      framework for ``derived stacks.''
      \index{derived algebraic geometry}
\end{enumerate}
\end{example}

\begin{remark}\label{rem:hott}
\index{homotopy type theory}
\index{univalence axiom}
$\infty$-toposes provide the semantic models for \textbf{homotopy type
theory} (HoTT).  The univalence axiom of Voevodsky holds in every
$\infty$-topos, and the internal language of an $\infty$-topos is
(conjecturally) homotopy type theory.  This is an active area of research
connecting category theory, homotopy theory, and the foundations of
mathematics.
\end{remark}

\index{infinity-topos@$\infty$-topos|)}


%% ============================================================
\section{Exercises}\label{sec:infty-exercises}
%% ============================================================

\begin{exercise}\label{ex:nerve-quasi}
\index{nerve!inner horn condition}
Let $\mathcal{C}$ be an ordinary category.  Show that the nerve
$N(\mathcal{C})$ is a quasi-category and that inner horn fillers
are unique.  Deduce that $N$ is fully faithful as a functor
$\Cat\to\mathrm{sSet}$.
\end{exercise}

\begin{exercise}\label{ex:kan-groupoid}
\index{Kan complex!groupoid property}
Show that a quasi-category~$X$ is a Kan complex if and only if every
morphism (1-simplex) in~$X$ is an equivalence (invertible up to a
2-simplex).
\end{exercise}

\begin{exercise}\label{ex:homotopy-cat-nerve}
\index{homotopy category!of a nerve}
Let $\mathcal{C}$ be an ordinary category.  Show that
$\mathrm{h}(N(\mathcal{C})) \cong \mathcal{C}$.
\end{exercise}

\begin{exercise}\label{ex:infty-presheaf-colim}
\index{infinity-presheaf@$\infty$-presheaf!as colimit of representables}
Let $\mathcal{C}$ be a small $(\infty,1)$-category and
$F\in\mathcal{P}(\mathcal{C})$ an $\infty$-presheaf.  Using the
$\infty$-Yoneda lemma, show that~$F$ is a colimit of representable
presheaves.
\end{exercise}

\begin{exercise}\label{ex:infty-adj-unique}
Show that if $F\colon\mathcal{C}\to\mathcal{D}$ is a functor of
$(\infty,1)$-categories that admits a right adjoint~$G$, then $G$ is
unique up to a contractible space of choices.
\textit{Hint:} use the $\infty$-Yoneda lemma to express $G(D)$ via
the presheaf $C\mapsto\mathrm{Map}_\mathcal{D}(FC, D)$.
\end{exercise}

\begin{exercise}\label{ex:infty-topos-slice}
\index{infinity-topos@$\infty$-topos!slice}
Let $\mathcal{E}$ be an $\infty$-topos and $X\in\mathcal{E}$ an object.
Show that the slice $\infty$-category $\mathcal{E}_{/X}$ is again an
$\infty$-topos.
\end{exercise}

\begin{exercise}[Descent]\label{ex:infty-descent}
\index{descent}
\index{infinity-topos@$\infty$-topos!descent}
Let $\mathcal{E}$ be an $\infty$-topos and $f\colon X\to Y$ a morphism
in~$\mathcal{E}$.  Show that the pullback functor
$f^*\colon\mathcal{E}_{/Y}\to\mathcal{E}_{/X}$ preserves all colimits.
This is the $\infty$-categorical formulation of \textbf{descent}.
\end{exercise}

\index{infinity-category@$\infty$-category|)}
\index{higher category theory|)}


%%%%%%%%%%%%%%%%%%%%%%%%%%%%%%%%%%%%%%%%%%%%%%%%%%%%%%%%%%%%
%%                  END MATTER                            %%
%%%%%%%%%%%%%%%%%%%%%%%%%%%%%%%%%%%%%%%%%%%%%%%%%%%%%%%%%%%%

\appendix

\chapter{Review of Algebra and Topology}\label{app:review}
\minitoc

\vspace{1em}

We collect here the basic definitions from algebra and topology that
are used throughout the text.  The reader familiar with these notions
may safely skip this appendix and refer back to it as needed.

%% ============================================================
\section{Groups, rings, and modules}\label{sec:app-algebra}
%% ============================================================

\begin{definition}[Group]\label{def:app-group}
\index{group!definition}
A \textbf{group} is a set~$G$ equipped with a binary operation
$\cdot\colon G\times G\to G$, an identity element $e\in G$, and an
inverse map $(-)^{-1}\colon G\to G$, satisfying associativity,
identity, and inverse laws.  A group is \textbf{abelian} if
$a\cdot b = b\cdot a$ for all $a,b\in G$.
\end{definition}

\begin{definition}[Ring]\label{def:app-ring}
\index{ring!definition}
A \textbf{ring} $(R, +, \cdot, 0, 1)$ is an abelian group $(R,+,0)$
equipped with an associative multiplication with unit~$1$, such that
multiplication distributes over addition.  A ring is \textbf{commutative}
if $ab = ba$ for all $a,b\in R$.
\end{definition}

\begin{definition}[Module]\label{def:app-module}
\index{module!definition}
Let $R$ be a ring.  A \textbf{left $R$-module} is an abelian group $(M,+)$
equipped with a scalar multiplication $R\times M\to M$ satisfying the
usual axioms: $r(m+n) = rm+rn$, $(r+s)m = rm+sm$, $(rs)m = r(sm)$,
$1m = m$.
\end{definition}

\begin{definition}[Field]\label{def:app-field}
\index{field!definition}
A \textbf{field} is a commutative ring~$\K$ in which every non-zero
element is invertible, i.e.\ $\K^\times = \K\setminus\{0\}$.
\end{definition}

\begin{definition}[Algebra over a field]\label{def:app-algebra}
\index{algebra!over a field}
A \textbf{$\K$-algebra} is a ring~$A$ that is also a $\K$-vector space,
with multiplication $\K$-bilinear.  Equivalently, it is a monoid object
in $(\Vect_\K, \otimes_\K, \K)$.
\end{definition}


%% ============================================================
\section{Topological spaces}\label{sec:app-topology}
%% ============================================================

\begin{definition}[Topological space]\label{def:app-top-space}
\index{topological space!definition}
A \textbf{topological space} is a pair $(X, \tau)$ where $X$ is a set
and $\tau\subseteq\mathcal{P}(X)$ is a collection of subsets (called
\emph{open sets}) satisfying: (i)~$\varnothing, X\in\tau$;
(ii)~$\tau$ is closed under arbitrary unions; (iii)~$\tau$ is closed
under finite intersections.
\end{definition}

\begin{definition}[Continuous map]\label{def:app-continuous}
\index{continuous map}
A map $f\colon X\to Y$ between topological spaces is \textbf{continuous}
if $f^{-1}(U)\in\tau_X$ for every $U\in\tau_Y$.
\end{definition}

\begin{definition}[Hausdorff space]\label{def:app-hausdorff}
\index{Hausdorff space}
A topological space~$X$ is \textbf{Hausdorff} (or $T_2$) if for every
pair of distinct points $x\ne y$ there exist disjoint open sets
$U\ni x$ and $V\ni y$.
\end{definition}

\begin{definition}[Compact space]\label{def:app-compact}
\index{compact space}
A topological space~$X$ is \textbf{compact} if every open cover has a
finite subcover.
\end{definition}

\begin{definition}[Connected space]\label{def:app-connected}
\index{connected space}
A topological space~$X$ is \textbf{connected} if it cannot be written
as a disjoint union of two non-empty open sets.
\end{definition}


%% ============================================================
\section{Simplicial sets}\label{sec:app-simplicial}
%% ============================================================

\begin{definition}[Simplex category]\label{def:app-simplex-cat}
\index{simplex category!definition}
The \textbf{simplex category}~$\Delta$ has as objects the finite
ordinals $[n] = \{0, 1, \ldots, n\}$ for $n\ge 0$, and as morphisms
the order-preserving maps.
\end{definition}

\begin{definition}[Simplicial set]\label{def:app-sset}
\index{simplicial set!definition}
A \textbf{simplicial set} is a functor $X\colon\Delta^{\op}\to\Set$.
It is specified by sets $X_n$ (the $n$-simplices) together with
face maps $d_i\colon X_n\to X_{n-1}$ and degeneracy maps
$s_i\colon X_n\to X_{n+1}$ satisfying the simplicial identities.
\end{definition}

\begin{definition}[Geometric realisation]\label{def:app-realisation}
\index{geometric realisation}
The \textbf{geometric realisation} of a simplicial set~$X$ is
\[
    |X| = \left(\coprod_{n\ge 0} X_n\times\Delta^n_{\mathrm{top}}\right)
    \Big/ {\sim}
\]
where $\Delta^n_{\mathrm{top}}$ is the topological $n$-simplex
and the equivalence relation is generated by the face and degeneracy maps.
\end{definition}


%% ============================================================
\section{Basic homological algebra}\label{sec:app-homological}
%% ============================================================

\begin{definition}[Chain complex]\label{def:app-chain-complex}
\index{chain complex!definition}
A \textbf{chain complex} over a ring~$R$ is a sequence of $R$-modules
and homomorphisms
\[
    \cdots \xrightarrow{d_{n+1}} C_n \xrightarrow{d_n} C_{n-1}
    \xrightarrow{d_{n-1}} \cdots
\]
with $d_n\circ d_{n+1} = 0$ for all~$n$.
The \textbf{homology} is $H_n(C_\bullet) = \ker d_n / \mathrm{im}\, d_{n+1}$.
\index{homology}
\end{definition}

\begin{definition}[Exact sequence]\label{def:app-exact}
\index{exact sequence}
A chain complex is \textbf{exact} at~$C_n$ if $H_n = 0$, i.e.\
$\ker d_n = \mathrm{im}\, d_{n+1}$.  A \textbf{short exact sequence}
is an exact sequence $0\to A\to B\to C\to 0$.
\end{definition}

\begin{definition}[Tensor product of modules]\label{def:app-tensor}
\index{tensor product!of modules}
For a commutative ring~$R$ and $R$-modules $M,N$, the \textbf{tensor
product} $M\otimes_R N$ is the $R$-module generated by symbols
$m\otimes n$ ($m\in M$, $n\in N$) subject to bilinearity relations.
It is characterised by the universal property:
$\Hom_R(M\otimes_R N, P) \cong \mathrm{Bilin}_R(M\times N, P)$.
\end{definition}

\begin{definition}[Projective and injective modules]\label{def:app-proj-inj}
\index{projective module}
\index{injective module}
An $R$-module~$P$ is \textbf{projective} if $\Hom_R(P,-)$ is exact.
An $R$-module~$Q$ is \textbf{injective} if $\Hom_R(-,Q)$ is exact.
\end{definition}


%% ============================================================
\section{Posets and lattices}\label{sec:app-posets}
%% ============================================================

\begin{definition}[Partially ordered set]\label{def:app-poset}
\index{partially ordered set}
\index{poset|see{partially ordered set}}
A \textbf{partially ordered set} (or \textbf{poset}) is a set~$P$
equipped with a relation~$\le$ that is reflexive, antisymmetric, and
transitive.  A poset may be viewed as a category with objects~$P$
and a unique morphism $x\to y$ whenever $x\le y$.
\end{definition}

\begin{definition}[Lattice]\label{def:app-lattice}
\index{lattice}
A \textbf{lattice} is a poset in which every pair of elements $\{a,b\}$
has a meet $a\wedge b$ (greatest lower bound) and a join $a\vee b$
(least upper bound).  A \textbf{complete lattice} has meets and joins
for all subsets.
\index{complete lattice}
\end{definition}

\begin{definition}[Heyting algebra]\label{def:app-heyting}
\index{Heyting algebra!definition}
A \textbf{Heyting algebra} is a lattice~$H$ with a binary operation
$\Rightarrow$ (implication) satisfying: $a\wedge b\le c$ if and only
if $a\le (b\Rightarrow c)$.  Every Boolean algebra is a Heyting algebra,
but not conversely.
\end{definition}


%% ============================================================
%% INDEX AND CLOSING
%% ============================================================

\cleardoublepage
\phantomsection
\addcontentsline{toc}{chapter}{Index}
\printindex

\end{document}
